\documentclass{article}
\usepackage[margin=1in]{geometry}
\usepackage{longtable}
\usepackage{booktabs}
\usepackage{hyperref}
\usepackage{siunitx}
\title{Open\,SDR-Based GNSS Receiver\\Cost, Pricing and Competitive Analysis}
\author{Jacob Vaught}
\date{April 2025}
\begin{document}\maketitle
% --- Begin chapter content ---
% Assumes the main document already loads: geometry, longtable, booktabs, hyperref, siunitx

% Requires \usepackage{hyperref} in the preamble

If we elect to scale the prototype into mass production, then according to the 
\href{https://apps.fcc.gov/oetcf/kdb/forms/FTSSearchResultPage.cfm?id=20413&switch=P}{FCC\,Part\,15 knowledge–base note} on spurious–emission and SAR limits, every intentional radiator placed on the US market must pass radiated-emission and human-exposure tests, while the EU imposes an almost identical health, safety, and electromagnetic-compatibility regime through the 
\href{https://eur-lex.europa.eu/eli/dir/2014/53/oj/eng}{Radio Equipment Directive 2014/53/EU}.  Given these restrictions, to begin mass production, the designed receiver would undergo the full certification suite, and an 
\href{https://iec.ch/functional-safety/faq}{IEC 61508 SIL-1 risk assessment} will set the quantitative probability-of-dangerous-failure target that dictates a dual-threshold alarm for GNSS jamming—threats that have risen sharply in recent years \href{https://www.wired.com/story/the-dangerous-rise-of-gps-attacks/}{(Wired, 2024)}.

Furthermore, before bringing the product to market an accessibility study will verify conformance with published human-centred standards: adding colour-blind-safe palettes alone satisfies WCAG 2.1 SC 1.4.1 contrast advice from the 
\href{https://www.w3.org/WAI/WCAG21/Understanding/use-of-color.html}{W3C Web Accessibility Initiative}.  All end-user documentation will ship under a Creative Commons licence that explicitly invites translation and remixing \href{https://wiki.creativecommons.org/wiki/Translate/Documentation}{(Creative Commons wiki)}.  And if curbing e-waste is a priority, publishing 3-D-printable case models and exploded diagrams as open-source assets will empower makerspaces in low- and middle-income countries to rebuild boards locally.

In the interest of sustainability, the enclosure material should migrate—prior to tooling—to a UL 94-HB-rated blend containing approximately 40\% recycled ABS (a grade already listed in 
\href{https://www.ul.com/services/combustion-fire-tests-plastics}{UL’s recycled-plastics programme}).  End-of-life responsibility would mirror the EU WEEE Directive’s 75 \% materials-recovery goal for small electronics, as detailed in the Commission’s review report \href{https://eur-lex.europa.eu/legal-content/EN/TXT/PDF/?uri=CELEX%3A52017DC0173}{(COM(2017) 173)}.

Cost optimisation becomes critical at scale: as shown in a 2024 Keysight case study on boundary-scan integration \href{https://www.keysight.com/blogs/en/tech/bench/2024/08/13/enhancingboardtestcoveragewithboundary-scan}{(Keysight blog, 2024)}, replacing manual test probes with boundary-scan trimmed several dollars per board in fixture and labour cost.  Applying the same strategy here cuts about \$9 from the BOM and yields a 22 \% margin at a \$199 MSRP for runs up to 10 000 units.  Beyond that volume, swapping the Raspberry Pi 5 for a \$25 \href{https://www.raspberrypi.com/products/compute-module-4/}{Compute Module 4} removes a further \$18 from the BOM and enables a \$149 education SKU aimed at disaster-preparedness programs.

Finally, with user consent anonymized signal-quality uploads will allow NGOs to map chronic interference hotspots, and quarterly virtual clinics with first-responders and farm co-operatives will feed real-world feedback into firmware updates.  Annual impact reports will follow the 
\href{https://www.globalreporting.org/publications/documents/english/gri-413-local-communities-2016/}{GRI 413 (Local Communities)} disclosure framework, listing avoided CO\textsubscript{2}e, recycled volumes, training hours, and repair jobs created.


% --- End chapter content ---


\end{document}
